\chapter{Macro systématique}\label{ch:s2:systematic-macro}

Value, momentum et carry, mesurés pays par pays dans les obligations, les devises et les indices boursiers, composent un portefeuille macro qu'aucun économiste n'a conçu.
Asness, Moskowitz et Pedersen ont mis en évidence des primes de value et de momentum sur huit marchés et classes d'actifs. Chaque prime était plus corrélée à ses homologues dans d'autres
classes d'actifs que les classes ne l'étaient entre elles, et value et momentum étaient corrélés négativement. Sur le marché synthétique de ce chapitre, formé de
trente contrats à terme, value affiche un ratio de Sharpe de 0.84, momentum 0.78, carry 0.33 et un \omterm{def:s2:systematic-macro:surprise}{signal de surprises économiques} 0.29. Presque décorrélés, ils se combinent
pour atteindre 1.15. Mesurer value à partir des prix plutôt qu'à partir d'un ancrage fondamental en fait l'image miroir du momentum, $-0.84$, et la combinaison tombe à
0.31. La construction est \texttt{firm.sysmacro}.

\section{Signaux sur les pays et les classes d'actifs}

L'univers est celui des contrats à terme synthétiques du livre~8 (chapitre~19)~: dix indices boursiers, dix obligations d'État et dix devises, un de chaque par pays, sur
trente ans, avec une dérive persistante (la tendance) et une prime de carry. \texttt{firm.sysmacro} ajoute deux effets implantés. Chaque marché porte un
écart de valorisation par rapport à la fair value, d'écart-type 10~\% et de demi-vie de trois ans~; le prix intègre cet écart, et sa correction est la
prime de value. Chaque pays dispose d'un indice de surprises économiques, de demi-vie 60 jours, qui fait bouger ses marchés dans les jours suivants.

\begin{definition}[Signal value macro]\label{def:s2:systematic-macro:value}
Un \emph{signal value macro}\index{signal value macro} classe les actifs des pays selon leur prix par rapport à un ancrage fondamental, comme la parité de
pouvoir d'achat pour une devise, un modèle de rendement réel ou de prime de terme pour une obligation, ou les bénéfices et la valeur comptable pour un indice boursier, en achetant les moins chers et en vendant les
plus chers.
\end{definition}

Chaque signal est négocié en long-short au sein de chaque classe d'actifs, réajusté chaque mois selon le rang, et calibré à une volatilité de 10~\% sur son trimestre glissant
(\cref{lst:s2:systematic-macro:signals}). Value repose sur un ancrage que l'analyste de chaque marché se trompe de 10~\% de façon persistante. Le momentum est le
rendement sur douze mois hors du dernier mois (livre~8, chapitre~19), et le carry est celui du livre~8 (chapitre~20). La variante de value fondée sur les prix est l'opposé
de la variation au comptant des cinq dernières années, soit le rendement diminué du carry perçu.

\section{Signaux de données macroéconomiques}

\begin{definition}[Momentum macro]\label{def:s2:systematic-macro:momentum}
Le \emph{momentum macro}\index{momentum macro} prend des positions longues sur les actifs des pays dont les tendances économiques fondamentales (croissance, inflation,
balance extérieure, politique monétaire) s'améliorent et prend des positions courtes là où elles se dégradent, le pendant du momentum sur les prix.
\end{definition}

\begin{definition}[Signal de surprises économiques]\label{def:s2:systematic-macro:surprise}
Un \emph{signal de surprises économiques}\index{signal de surprises économiques} mesure de combien les dernières publications de données d'un pays ont dépassé ou manqué
des prévisions, et positionne ses actifs en vue de la réaction différée des marchés~: actions et devise à la hausse, obligations à la baisse après des surprises positives.
\end{definition}

Brooks décrit le \omterm{def:s2:systematic-macro:momentum}{momentum macro} comme une stratégie macro globale systématique, longue sur les actifs dont les tendances macroéconomiques s'améliorent et courte sur ceux
dont les tendances se dégradent, avec une faible corrélation aux classes d'actifs traditionnelles et une diversification dans les marchés baissiers d'actions et les environnements à
rendements en hausse. Le \omterm{def:s2:systematic-macro:surprise}{signal de surprises économiques} synthétique en est une version simple. L'indice de chaque pays fait monter son indice boursier de 4~\% par an et par unité, sa
devise de 3~\%, et baisser son obligation de 3~\%, et le trader voit l'indice avec un bruit égal à la moitié de sa taille.

\section{Combiner les signaux}

\begin{center}
\small
\begin{tabular}{@{}lccc@{}}
\toprule
24.5 ans, 10~\% de volatilité chacune & Sharpe & drawdown max. & corr.\ avec le momentum\\
\midrule
value (ancrage) & 0.84 & $-25~\%$ & $-0.03$\\
value à partir des prix (retournement) & $-0.84$ & & $-0.43$\\
momentum & 0.78 & $-23~\%$ & 1\\
carry & 0.33 & $-32~\%$ & 0.04\\
surprises & 0.29 & $-38~\%$ & $-0.01$\\
\midrule
les quatre & 1.15 & $-21~\%$ &\\
les quatre, value à partir des prix & 0.31 & &\\
tous sauf value & 0.81 & &\\
\bottomrule
\end{tabular}
\end{center}

Quatre signaux presque décorrélés se combinent comme l'arithmétique le prévoit. Des budgets de risque égaux donnent un ratio de Sharpe voisin de la somme des quatre divisée par la
racine carrée de quatre, 1.12, contre 1.15 mesuré (\cref{fig:s2:systematic-macro:cumulative}). Le pire drawdown de la combinaison, 21~\%, est inférieur à
celui de n'importe quel signal pris seul. Value mérite sa place~: sans lui, la combinaison tombe à 0.81.

La manière de mesurer value décide de son utilité. Sur ce marché, la tendance dure environ un an, et la variation de cours des cinq dernières années en porte encore une
partie. Prendre l'opposé de cette variation revient à être court sur les marchés qui ont progressé, si bien que le signal value fondé sur les prix devient le miroir du momentum~: $-0.84$ contre 0.78, une
corrélation de $-0.43$. Combinés, les deux s'annulent, et le portefeuille à quatre signaux ne rapporte que 0.31. Asness et ses coauteurs ont constaté que value et momentum étaient tous deux
positifs et corrélés négativement. Un marché synthétique ne reproduit cela que si l'ancrage de value n'est pas bâti sur les mêmes prix que ceux qu'utilise le momentum.

\begin{omfigure}
\begin{tikzpicture}
\begin{axis}[width=11cm,height=5.6cm,xlabel={\small année},ylabel={\small rendement cumulé (\%, somme simple)},tick label style={font=\footnotesize},
  xmin=0,xmax=24.5,ymin=-20,ymax=300,grid=major,grid style={gray!25},table/col sep=comma,
  legend style={font=\scriptsize,at={(0.5,-0.3)},anchor=north},legend columns=5]
\addplot[omDef,thick] table[x=year,y=value]{figdata/strategies-2/16-systematic-macro/cumulative.csv};
\addplot[omThm,thick] table[x=year,y=momentum]{figdata/strategies-2/16-systematic-macro/cumulative.csv};
\addplot[gray,thick] table[x=year,y=carry]{figdata/strategies-2/16-systematic-macro/cumulative.csv};
\addplot[gray!60,thick,dashed] table[x=year,y=surprise]{figdata/strategies-2/16-systematic-macro/cumulative.csv};
\addplot[black,very thick] table[x=year,y=combined]{figdata/strategies-2/16-systematic-macro/cumulative.csv};
\legend{value,momentum,carry,surprises,combiné}
\end{axis}
\end{tikzpicture}

\omcaption{Quatre signaux macro synthétiques, chacun calibré à une volatilité de 10~\% et négocié en long-short au sein des actions, des obligations et des devises, ainsi que leur
combinaison à risque égal (également à 10~\%), sommes cumulées des rendements quotidiens sur 24.5 ans. Données~: \texttt{s2\_sysmacro.books}.}\label{fig:s2:systematic-macro:cumulative}
\end{omfigure}

\section{Ce que montre la littérature}

Le bilan public confirme la configuration, non les chiffres. Asness et ses coauteurs ont trouvé des primes de value et de momentum partout où ils ont regardé, avec une structure factorielle commune
entre les classes d'actifs et un lien partiel avec la liquidité mondiale de financement. Le \omterm{def:s2:systematic-macro:momentum}{momentum macro} de Brooks ajoute les tendances des données économiques. Les
ratios de Sharpe du chapitre sont ceux d'un marché à effets implantés (\cref{fig:s2:systematic-macro:sharpe}). Ils servent à montrer comment les signaux se combinent, et non
à prévoir ce que rapporterait un portefeuille réel. Ce qui se transpose est la méthode~: mesurer chaque signal isolément, mesurer leurs corrélations, attribuer à chacun un budget de risque
et vérifier que value et momentum ne sont pas mesurés à partir des mêmes données.

\begin{omfigure}
\begin{tikzpicture}
\begin{axis}[width=11.5cm,height=5.6cm,ybar,bar width=12pt,ylabel={\small ratio de Sharpe},tick label style={font=\footnotesize},
  xmin=0.4,xmax=8.6,ymin=-1.3,ymax=1.4,xtick={1,2,3,4,5,6,7,8},
  xticklabels={value,retournement,momentum,carry,surprises,{les quatre},{les quatre, value sur les prix},{tous sauf value}},
  x tick label style={font=\scriptsize,rotate=25,anchor=north east},
  grid=major,grid style={gray!25},table/col sep=comma,nodes near coords,every node near coord/.append style={font=\scriptsize}]
\addplot[fill=omThm!70,draw=omThm] table[x=x,y=sr]{figdata/strategies-2/16-systematic-macro/sharpe.csv};
\end{axis}
\end{tikzpicture}

\omcaption{Ratios de Sharpe des signaux macro synthétiques et de trois combinaisons~: les quatre (value à partir de l'ancrage), les quatre avec value à partir des prix,
et trois sans value. Données~: \texttt{s2\_sysmacro.stats}.}\label{fig:s2:systematic-macro:sharpe}
\end{omfigure}

\section{Fiches de stratégie}

\begin{strategyfile}{Value entre pays}\label{strat:s2:systematic-macro:value}
\sfield{Qui rémunère la stratégie, et pourquoi} Des investisseurs qui extrapolent, et un ancrage fondamental vers lequel les prix reviennent sur plusieurs années.
\sfield{Instruments et plateformes de négociation} Contrats à terme sur indices actions, sur obligations et sur devises, et forwards.
\sfield{Signal} Prix rapporté à un ancrage fondamental~: parité de pouvoir d'achat, rendements réels, rendements des bénéfices.
\sfield{Dimensionnement et exécution} Long-short au sein de chaque classe d'actifs~; rééquilibrage mensuel~; budgets de risque.
\sfield{Coûts} Rolls des contrats à terme~; faible rotation.
\sfield{Comment elle meurt} Ancrages faux pendant des années~; tendances qui vont plus loin que la value ne peut attendre.
\sfield{Horizon, capacité, infrastructure} Années~; grande capacité.
\sfield{Backtest honnête} Ancrages tels qu'ils pouvaient être calculés à l'époque, et non des données révisées.
\sfield{Sources} Asness, Moskowitz et Pedersen (2013)~; ce chapitre~: 0.84 à partir d'un ancrage, $-0.84$ à partir des prix.
\end{strategyfile}

\begin{strategyfile}{Momentum macro}\label{strat:s2:systematic-macro:momentum}
\sfield{Qui rémunère la stratégie, et pourquoi} Des marchés qui intègrent lentement les nouvelles économiques dans les prix.
\sfield{Instruments et plateformes de négociation} Comme ci-dessus.
\sfield{Signal} Tendances de la croissance, de l'inflation, des soldes extérieurs et de la politique monétaire~; momentum des prix en complément.
\sfield{Dimensionnement et exécution} Long-short~; mensuel.
\sfield{Coûts} Rotation modérée.
\sfield{Comment elle meurt} Points de retournement~; révisions des données.
\sfield{Horizon, capacité, infrastructure} Mois~; données macro point-in-time.
\sfield{Backtest honnête} Premières publications des données, non révisées.
\sfield{Sources} Brooks (2017)~; ce chapitre~: momentum des prix 0.78.
\end{strategyfile}

\begin{strategyfile}{Carry entre classes d'actifs}\label{strat:s2:systematic-macro:carry}
\sfield{Qui rémunère la stratégie, et pourquoi} Comme dans le livre~8, chapitre~20~: détenteurs d'actifs à faible carry, et risque de krach pesant sur les carry élevés.
\sfield{Instruments et plateformes de négociation} Contrats à terme et forwards dans trois classes d'actifs.
\sfield{Signal} Carry de chaque marché.
\sfield{Dimensionnement et exécution} Long-short au sein de chaque classe~; budgets de risque.
\sfield{Coûts} Rolls.
\sfield{Comment elle meurt} Krachs du carry.
\sfield{Horizon, capacité, infrastructure} Mois.
\sfield{Backtest honnête} Carry calculé à partir des courbes de contrats à terme, et non de taux supposés.
\sfield{Sources} Livre~8, chapitre~20~; ce chapitre~: 0.33.
\end{strategyfile}

\begin{strategyfile}{Biais de surprises économiques}\label{strat:s2:systematic-macro:surprise}
\sfield{Qui rémunère la stratégie, et pourquoi} Des investisseurs qui réagissent lentement aux surprises statistiques.
\sfield{Instruments et plateformes de négociation} Contrats à terme sur indices actions, obligations et devises de chaque pays.
\sfield{Signal} Dernières publications de données rapportées aux prévisions consensuelles, pays par pays.
\sfield{Dimensionnement et exécution} Long-short~; rééquilibrage à mesure de l'arrivée des données.
\sfield{Coûts} Rotation plus élevée que celle de la value.
\sfield{Comment elle meurt} Les prévisionnistes et les marchés rattrapent le retard plus vite.
\sfield{Horizon, capacité, infrastructure} Semaines à mois~; un calendrier de données assorti de prévisions.
\sfield{Backtest honnête} Heures de publication et premières estimations.
\sfield{Sources} Ce chapitre~: 0.29~; aucun chiffre de performance public vérifié.
\end{strategyfile}

\begin{strategyfile}{Portefeuille combiné de facteurs macro}\label{strat:s2:systematic-macro:combined}
\sfield{Qui rémunère la stratégie, et pourquoi} La somme des stratégies ci-dessus.
\sfield{Instruments et plateformes de négociation} Ceux de toutes les stratégies ci-dessus.
\sfield{Signal} Value, momentum, carry et données macro, avec des budgets de risque égaux ou estimés.
\sfield{Dimensionnement et exécution} Volatilité cible au niveau du portefeuille.
\sfield{Coûts} La compensation des positions réduit la rotation.
\sfield{Comment elle meurt} Corrélations qui montent ensemble en période de crise.
\sfield{Horizon, capacité, infrastructure} Mois~; un portefeuille de contrats à terme multi-actifs.
\sfield{Backtest honnête} Signaux définis avant le test~; corrélations mesurées hors échantillon.
\sfield{Sources} Asness, Moskowitz et Pedersen (2013)~; ce chapitre~: 1.15 combiné.
\end{strategyfile}

\section{Tutoriel~: conçu par personne}\label{tut:s2:systematic-macro:tut}

\begin{tutorial}
\textbf{Objectif.} Introduire des effets de value et de surprises macro dans les contrats à terme synthétiques, construire cinq signaux et les combiner. \textbf{État final~:} le tableau
et les deux figures.
\begin{tutsteps}
\item \textbf{Signaux et portefeuilles}.
\omcode{code/firm/sysmacro/firm_sysmacro.py}{75}{114}{Les cinq signaux, les pondérations long-short au sein de chaque classe et le portefeuille à volatilité cible.}\label{lst:s2:systematic-macro:signals}
\item \textbf{Les combinaisons}.
\omcode{code/strategies-2/16-systematic-macro/python/s2_sysmacro.py}{30}{51}{Chaque signal, trois combinaisons et leurs statistiques.}\label{lst:s2:systematic-macro:books}
\item \textbf{Lancer} \texttt{stats()}, \texttt{correlations()} et \texttt{fig\_sysmacro.py}.
\end{tutsteps}
\textbf{Pour aller plus loin.} Estimer les budgets de risque à partir des corrélations glissantes au lieu de budgets égaux~; ajouter une erreur d'ancrage lente qui dérive~;
tester la value fondée sur les prix sur un marché dont la tendance dure trois mois.
\end{tutorial}

\section{Construction~: macro systématique}\label{bld:s2:systematic-macro:sysmacro}

\begin{build}
\bfield{Rôle} Signaux de value, de momentum, de carry et de surprises macro transversaux aux classes d'actifs, et leur combinaison à budgets de risque.
\bfield{Interface} \texttt{MacroConfig(\dots)}, \texttt{macro\_market(cfg)}, \texttt{signal(m, name, t)}, \texttt{portfolio(m, name, cfg)},
\texttt{combine(books, cfg)}.
\bfield{Règles} Signaux issus des données jusqu'à la veille~; long-short au sein de chaque classe~; chaque portefeuille dimensionné sur sa propre volatilité glissante.
\bfield{Tests d'acceptation} \texttt{code/firm/sysmacro/tests/}~: les pondérations au sein des classes sont neutres avec un brut unitaire~; les formes de marché et les signaux~; la
combinaison vise sa volatilité cible.
\bfield{Pour aller plus loin} Budgets de risque estimés~; données macro réelles~; coûts de transaction.
\end{build}

\begin{omsources}
\item C.~S.~Asness, T.~J.~Moskowitz et L.~H.~Pedersen, ``Value and momentum everywhere'', \emph{Journal of Finance} 68(3), 2013.
\item J.~Brooks, ``A half century of macro momentum'', livre blanc AQR, 2017.
\end{omsources}

\section{Exercices}

\begin{exercise}[$\star$]\label{exo:s2:systematic-macro:1}
Quatre signaux non corrélés ont des ratios de Sharpe de 0.84, 0.78, 0.33 et 0.29, et des budgets de risque égaux. Calculer le ratio de Sharpe de leur combinaison.
\end{exercise}

\begin{exercise}[$\star$]\label{exo:s2:systematic-macro:2}
Deux signaux ont des ratios de Sharpe de 0.84 et 0.78 et une corrélation de $-0.43$. Calculer le ratio de Sharpe d'une combinaison à risque égal.
\end{exercise}

\begin{exercise}[$\star$]\label{exo:s2:systematic-macro:3}
Un écart de valorisation a une demi-vie de trois ans. Calculer la part de cet écart qui se referme en un an.
\end{exercise}

\begin{exercise}[$\star\star$]\label{exo:s2:systematic-macro:4}
Pourquoi la valeur mesurée à partir des variations de prix sur cinq ans devient-elle l'image miroir du momentum sur ce marché~?
\end{exercise}

\begin{exercise}[$\star\star$]\label{exo:s2:systematic-macro:5}
Pourquoi un backtest sur données macroéconomiques doit-il utiliser les données de première publication~?
\end{exercise}

\begin{exercise}[$\star\star$]\label{exo:s2:systematic-macro:6}
Pourquoi négocier les signaux en long-short au sein de chaque classe d'actifs plutôt qu'entre elles~?
\end{exercise}

\begin{exercise}[$\star\star\star$]\label{exo:s2:systematic-macro:7}
\emph{Programmation.} Relancer \texttt{macro\_market} avec \texttt{MacroConfig(anchor\_noise=0.3)}. Que se passe-t-il pour la valeur, et pour la combinaison~?
\end{exercise}

\begin{exercise}[$\star\star\star$]\label{exo:s2:systematic-macro:8}
\emph{Trouver la faille.} «~Notre combinaison affiche un ratio de Sharpe de 1.15, nous allons donc la faire tourner à une volatilité de 30~\%.~»
\end{exercise}

\section{Problème~: Conçu par personne}

\begin{problem}[{Problème du week-end --- un portefeuille macro systématique}]\label{pb:s2:systematic-macro:1}
Les contrats à terme synthétiques du chapitre et les archives publiques.

\medskip\noindent\textbf{Partie I --- Signaux.}
\begin{enumerate}
\item Définir le \omterm{def:s2:systematic-macro:value}{signal value macro} et donner trois ancres.
\item Définir le \omterm{def:s2:systematic-macro:momentum}{momentum macro} et le \omterm{def:s2:systematic-macro:surprise}{signal de surprises économiques}.
\item Comment chaque signal est-il négocié et dimensionné~?
\item Qu'ont établi Asness, Moskowitz et Pedersen~?
\end{enumerate}

\medskip\noindent\textbf{Partie II --- Le marché fabriqué.}
\begin{enumerate}[resume]
\item Décrire l'écart de valorisation et l'indice de surprise.
\item Comment l'analyste appréhende-t-il chacun~?
\item Pourquoi l'erreur sur l'ancre importe-t-elle~?
\item Que décrit Brooks~?
\end{enumerate}

\medskip\noindent\textbf{Partie III --- Résultats.}
\begin{enumerate}[resume]
\item Donner le ratio de Sharpe de chaque signal.
\item Donner les corrélations.
\item Donner les ratios de Sharpe des trois combinaisons.
\item Pourquoi la valeur calculée à partir des prix échoue-t-elle ici~?
\end{enumerate}

\medskip\noindent\textbf{Partie IV --- Le verdict.}
\begin{enumerate}[resume]
\item Énoncer le \emph{résultat nommé}~: le ratio de Sharpe de chaque signal et celui de la combinaison, avec les corrélations qui la rendent possible.
\item Que prédit l'arithmétique des signaux non corrélés~?
\item Que devient la combinaison en période de crise~?
\item Comment fixer les budgets de risque~?
\item Comment backtester honnêtement le portefeuille~?
\item Quelle fiche de stratégie affiche le turnover le plus élevé~?
\item Quel lien ce chapitre entretient-il avec les chapitres du livre~8 consacrés à la tendance et au carry~?
\item En une phrase~: pourquoi une combinaison de signaux simples fonctionne-t-elle~?
\end{enumerate}
\end{problem}

\section{Questions d'entretien}

\begin{interviewq}[$\star$ \iqroles{chercheur}]\label{iq:s2:systematic-macro:1}
Comment mesureriez-vous la valeur des devises, des obligations et des indices actions~?
\end{interviewq}

\begin{interviewq}[$\star\star$ \iqroles{chercheur}]\label{iq:s2:systematic-macro:2}
Pourquoi la valeur et le momentum sont-ils corrélés négativement, et pourquoi cela est-il utile~?
\end{interviewq}

\begin{interviewq}[$\star\star$ \iqroles{trader}]\label{iq:s2:systematic-macro:3}
Le signal de momentum de votre portefeuille macro est vendeur sur les obligations japonaises, et son signal de valeur est acheteur. Que faites-vous~?
\end{interviewq}

\begin{interviewq}[$\star\star$ \iqroles{risque}]\label{iq:s2:systematic-macro:4}
Comment soumettriez-vous un portefeuille macro systématique à des tests de résistance~?
\end{interviewq}

\begin{interviewq}[$\star\star$ \iqroles{développeur}]\label{iq:s2:systematic-macro:5}
Concevez une base de données macroéconomiques point-in-time pour les backtests.
\end{interviewq}

\begin{interviewq}[$\star\star\star$ \iqroles{chercheur}]\label{iq:s2:systematic-macro:6}
Montrez que les combinaisons à risque égal de $n$ signaux, de ratios de Sharpe $S_i$ et de corrélation commune deux à deux $\rho$, ont un ratio de Sharpe égal à $\sum_i S_i /
\sqrt{n + n(n - 1)\rho}$.
\end{interviewq}
